%% ma: L12-B01-0006
%% lop: 12
%% chuong: 1
%% bai: 1
%% dang: 12.1.5
%% loai: TN4
%% muc: TH
%% dap_an: "C"
%% nguon: "(NBV)-ĐỀ SỐ 5-MỖI NGÀY 1 ĐỀ THI 2025.pdf (D05, MỖI NGÀY 1 ĐỀ - Đề số 5), Phần I Câu 1"
%% kien_thuc: "Bài 1 (tính đơn điệu của hàm số; đạo hàm)"
%% trang_thai: cho-duyet
%% ly_do: "Dạng toán mới đề xuất 12.1.5: nhận biết hàm số nghịch biến trên R trong bốn hàm cho trước"
\begin{ex}
Trong các hàm số sau, hàm số nào nghịch biến trên $\mathbb R$?
\choice{$y=2^x$}{$y=\dfrac{x+3}{2x-2}$}{\True $y=-x^3+3x^2-3x+2$}{$y=\cos x$}
\loigiai{Với $y=-x^3+3x^2-3x+2$: $y'=-3x^2+6x-3=-3(x-1)^2\le0$ với mọi $x$, dấu bằng chỉ tại $x=1$, nên hàm số nghịch biến trên $\mathbb R$.
Các hàm còn lại không thỏa: $y=2^x$ đồng biến trên $\mathbb R$; $y=\dfrac{x+3}{2x-2}$ không xác định tại $x=1$ (chỉ nghịch biến trên từng khoảng); $y=\cos x$ không đơn điệu trên $\mathbb R$.}
\end{ex}
